\section{Анализ предметной области}

\subsection{Понятие фитнеса, здорового образа жизни и роль информационных систем}

Фитнес в контексте прикладной информатики следует рассматривать как совокупность методов физической активности, направленных на развитие силы, выносливости, гибкости, координации и общего функционального состояния организма. Здоровый образ жизни включает не только регулярные тренировки, но и контроль режима питания, восстановления, сна, уровня стрессовой нагрузки и соблюдение разумного баланса между физической активностью и отдыхом.

Для пользователя, который занимается самостоятельно, особенно важны регулярность занятий, наглядность прогресса и возможность адаптировать план под свои цели. Именно в этом контексте информационные системы приобретают практическую ценность: они позволяют автоматизировать учет тренировок, хранить историю результатов, формировать рекомендации на основе введенных данных и поддерживать мотивацию за счет визуализации динамики.

В рамках данной работы информационная система рассматривается как инструмент поддержки принятия решений пользователем. Она не заменяет тренера или врача, однако помогает:
\begin{itemize}
	\item формировать индивидуальный план тренировок на основе целей, уровня подготовки и параметров пользователя;
	\item хранить историю выполненных упражнений и отслеживать прогресс;
	\item анализировать динамику нагрузки, частоту занятий и результативность;
	\item учитывать питание и другие сопутствующие параметры в упрощенной форме;
	\item повышать вовлеченность пользователя за счет интерфейса поддержки и социальной активности.
\end{itemize}

\subsection{Специфика задачи индивидуального планирования тренировок}

Составление индивидуального плана тренировок является не шаблонной задачей, поскольку у разных пользователей различаются цели, исходная физическая форма, ограничение по времени, уровень подготовки и доступное оборудование. Для одного пользователя приоритетом может быть снижение массы тела, для другого -- набор мышечной массы, для третьего -- поддержание физической активности и улучшение самочувствия.

При ручном ведении тренировочного процесса пользователь часто сталкивается со следующими проблемами:
\begin{itemize}
	\item сложность подбора упражнений, соответствующих уровню подготовки;
	\item отсутствие системного подхода к нагрузке и восстановлению;
	\item трудности в оценке прогресса без накопленной статистики;
	\item пропуск занятий из-за отсутствия напоминаний и наглядной структуры;
	\item невозможность быстро скорректировать план при изменении цели или состояния.
\end{itemize}

Следовательно, информационная система должна не только хранить данные, но и поддерживать логику персонализации, то есть формировать план занятий с учетом индивидуальных характеристик пользователя и предоставлять инструменты для последующей корректировки этого плана.

\subsection{Анализ пользовательских потребностей}

Для разработки информационной системы важно определить, какие именно задачи пользователь хочет решать в процессе регулярных тренировок. Практика показывает, что для большинства людей недостаточно просто иметь список упражнений: требуется система, которая помогает планировать занятия, отслеживать выполнение и наглядно показывает результат. Наиболее значимыми требованиями являются персонализация, мониторинг прогресса, хранение истории и удобство доступа к данным.

\begin{figure}[h]
	\centering
	\includegraphics[width=0.9\linewidth]{fitness_analysis_factors.png}
	\caption{Ключевые требования к системе индивидуальных тренировок (оценка значимости)}
	\label{fig:fitness-factors}
\end{figure}

Из рисунка ~\ref{fig:fitness-factors} видно, что приоритет следует отдавать функциям, связанным с индивидуальным планированием и анализом прогресса. Именно эти возможности формируют практическую ценность системы для пользователя и отличают ее от простого справочника упражнений.

\subsection{Классификация цифровых решений в области физической активности}

Современные цифровые решения для фитнеса и контроля физической активности можно разделить на несколько основных групп.

\subsubsection{Тренировочные дневники и трекеры активности}

Данные приложения ориентированы на учет упражнений, подходов, повторений, веса отягощений, времени выполнения и частоты тренировок. Их сильная сторона заключается в удобном журналировании и последующем анализе динамики. Однако такие системы часто ограничиваются фиксацией результата и не всегда предлагают гибкую модель построения индивидуального плана.

\subsubsection{Сервисы контроля питания}

Приложения для учета питания позволяют отслеживать калорийность, белки, жиры и углеводы, а также сравнивать фактическое потребление с целевыми значениями. Они полезны при работе с задачами снижения или набора массы, но обычно слабо связаны с тренировочной частью и не формируют единую картину физической подготовки пользователя.

\subsubsection{Экосистемы носимых устройств}

Фитнес-браслеты, умные часы и связанные с ними мобильные приложения предоставляют данные о пульсе, шаговой активности, сне и иногда о тренировочной нагрузке. Такие решения удобны за счет автоматического сбора информации, однако их функциональность зависит от конкретного оборудования, а логика работы часто привязана к закрытой экосистеме производителя.

\subsubsection{Социальные и мотивационные платформы}

Некоторые цифровые продукты делают акцент на челленджах, рейтингах, сообществах и публичной демонстрации результатов. Это повышает вовлеченность, но может быть недостаточно полезно для точного планирования тренировочного процесса. Кроме того, чрезмерная ориентация на социальное сравнение не всегда способствует устойчивому результату.

Таким образом, на практике существует разрыв между приложениями для учета активности, дневниками питания и мотивационными платформами. Это обосновывает разработку системы, объединяющей несколько функций в едином интерфейсе.

\subsection{Сравнение архитектурных решений}

При проектировании информационной системы для составления индивидуального плана тренировок возможны различные архитектурные подходы.

Подход только клиент подходит для простых локальных решений, но не обеспечивает надежное хранение данных, совместную работу с несколькими устройствами и централизованное управление учетными записями. Такой вариант удобен для прототипов, но ограничен при дальнейшем развитии системы.

Подход клиент + сервер + база данных является наиболее рациональным для рассматриваемой задачи. Он позволяет разделить пользовательский интерфейс, бизнес-логику и хранение данных, обеспечить авторизацию, синхронизацию прогресса и расширяемость системы. Именно такой вариант соответствует целям преддипломной практики и дает возможность реализовать полноценный учебный программный продукт.

Микросервисная архитектура, хотя и является современной и масштабируемой, для учебного проекта избыточна. Она увеличивает сложность развертывания, сопровождения и тестирования, не давая существенных преимуществ на данном этапе разработки.

Наиболее рациональным вариантом является клиент-серверная архитектура с базой данных. Она обеспечивает разделение интерфейса, бизнес-логики и хранения информации, а также позволяет реализовать авторизацию, синхронизацию пользовательских данных и дальнейшее расширение функциональности.

\subsection{Пользовательские сценарии и функциональные потребности}

Основной пользователь системы -- человек, который хочет самостоятельно заниматься физической активностью и контролировать результат. Для такого пользователя наиболее типичны следующие сценарии:
\begin{itemize}
	\item регистрация и авторизация в системе;
	\item создание и редактирование профиля с индивидуальными параметрами;
	\item получение персонального плана тренировок;
	\item ведение дневника выполненных упражнений;
	\item просмотр истории, статистики и динамики результатов;
	\item использование дополнительных разделов, связанных с питанием, поддержкой и социальным взаимодействием.
\end{itemize}

Из этих сценариев следует, что система должна обеспечивать не только ввод и хранение данных, но и удобную навигацию, быстрый доступ к ключевым функциям и возможность последующего расширения без полной переработки архитектуры.

\subsection{Требования к качеству данных и безопасности}

Так как система работает с персональными данными пользователя, к ней предъявляются требования по защите информации. Минимально необходимыми являются разграничение доступа к учетной записи, безопасное хранение паролей в виде хэшей, защита пользовательских данных от несанкционированного просмотра и корректная работа механизма аутентификации.

С точки зрения качества данных важно обеспечить целостность и актуальность информации. Данные о тренировках и результатах должны сохраняться в структурированном виде, чтобы их можно было использовать для построения истории прогресса и дальнейшего анализа. Для этого целесообразно применять реляционную модель данных и хранить основные сущности в отдельных таблицах.

С точки зрения интерфейса система должна быть понятной, предсказуемой и не перегруженной лишними действиями. Для пользователя, который занимается спортом регулярно, особенно важны быстрое внесение данных, наглядность результата и отсутствие избыточной сложности.

\subsection{Обзор аналогов}

На рынке существует множество цифровых решений, связанных с фитнесом и питанием, однако большинство из них решает только одну часть задачи. Одни приложения ориентированы на учет тренировок, другие -- на дневник питания, третьи -- на социальную мотивацию или работу с носимыми устройствами. При этом редко встречается решение, в котором в одном продукте объединены индивидуальный план тренировок, мониторинг результата, вспомогательные мотивационные механизмы и удобное веб-управление. Сравнительная характеристика функций аналогов представлена в таблице ~\ref{tab:analog-comparison}.

\begin{xltabular}{\textwidth}{|p{3.2cm}|X|X|}
	\caption{Сопоставление коммерческих решений и разрабатываемой информационной системы\label{tab:analog-comparison}}\\ \hline
	\centrow Решение & \centrow Основной функционал & \centrow Замечание для сопоставления с целевой ИС \\ \hline
	\endfirsthead
	\continuecaption{Продолжение таблицы \ref{tab:analog-comparison}}
	\centrow Решение & \centrow Основной функционал & \centrow Замечание для сопоставления с целевой ИС \\ \hline
	\finishhead
	
	Strong & План и учёт силовых тренировок, история подходов, базовая аналитика & Нет учёта питания, социального раздела и работы через браузер; расширенные возможности частично за подписку \\ \hline
	Hevy & Аналогичный учёт тренировок с акцентом на программы и статистику & Питание и полноценный социальный модуль не являются ядром продукта; веб-клиент ограничен; freemium-модель \\ \hline
	MyFitnessPal & Дневник питания, калории и нутриенты, штрихкоды, интеграции & План силовых тренировок и связанная аналитика прогресса не являются центральной задачей продукта \\ \hline
	Strava & Кардионагрузки, GPS-треки, сегменты, соревновательная мотивация & Фокус на беге и велосипеде; силовой зал и рацион как единый контур не объединены \\ \hline
	JEFIT & Программы упражнений, учёт тренировок, есть веб-доступ & Питание и сеть как у целевой ИС представлены слабо; часть функций за подпиской \\ \hline
	Lose It! & Контроль питания, цели по калориям, поддержка привычек & Как и у других дневников питания, тренировочный контур неразвит \\ \hline
	Наша ПИС & План и дневник тренировок, учёт прогресса, раздел питания, социальные механизмы, веб-клиент в одной системе & Учебная ИС; развитие ограничено учебной постановкой, а не лакунами по видам функций \\ \hline
\end{xltabular}

Наиболее близкие по функциональности системы, как правило, либо слишком узкоспециализированы, либо требуют подписки, либо завязаны на конкретную платформу. Поэтому разработка учебной информационной системы является обоснованной: она позволяет реализовать ключевые функции в одном программном комплексе и продемонстрировать полный цикл проектирования.

\subsection{Актуальность разработки}

Актуальность разработки обусловлена ростом интереса к цифровым сервисам для здоровья и физической активности, а также потребностью пользователей в простых и понятных инструментах самостоятельного контроля тренировочного процесса. Современный пользователь ожидает, что система не только сохранит данные, но и поможет организовать занятия, увидеть прогресс и поддерживать мотивацию на протяжении длительного времени.

Для преддипломной практики разработка подобной системы представляет интерес с точки зрения программной инженерии, так как включает проектирование предметной области, выбор архитектуры, разработку пользовательского интерфейса, организацию хранения данных и реализацию механизмов авторизации и обработки пользовательской информации.

